\documentclass[11pt]{article}
\usepackage[margin=1in]{geometry}
\usepackage[T1]{fontenc}
\usepackage[utf8]{inputenc}
\usepackage{lmodern,amsmath,amssymb,amsthm,microtype}
\usepackage[colorlinks=true,allcolors=blue]{hyperref}
\newtheorem{theorem}{Theorem}
\newtheorem{lemma}[theorem]{Lemma}
\newtheorem{proposition}[theorem]{Proposition}
\newcommand{\N}{\mathbb N}
\newcommand{\Z}{\mathbb Z}
\newcommand{\reach}{\textsc{Reachable}}
\newcommand{\nh}{\textsc{Nonhalting}}
\title{CoNP-completeness of binary directed chip-firing reachability}
\author{Alec Kriebel\\Independent researcher\\
\small ORCID: \href{https://orcid.org/0009-0001-9320-500X}{0009-0001-9320-500X}}
\date{15 October 2026, version 1.0}
\begin{document}
\maketitle
\begin{abstract}
Ordinary legal chip-firing reachability is coNP-complete under deterministic
polynomial-time many-one reductions for finite loopless strongly connected
directed multigraphs, with binary adjacency entries and nonnegative binary
initial and target configurations. We give an explicit reduction from the
complement of the halting problem of Farrell and Levine. A finite-state
cutoff controls an exact simulation, a gate supplies a cleanup reserve, and
a marker forces that gate to have fired in every target-reaching word.
Together with the prior coNP-membership theorem of Hujter, Kiss and
T\'othm\'er\'esz, this answers the historical general-digraph hardness
question in the stated model. We credit earlier period, recurrence and
reserve-simulation arguments. Exact-result, gadget and method priority,
independence of discovery, and present-day openness remain unknown after a
bounded source audit. This extensively AI-assisted note is unrefereed.
\end{abstract}

\section{The language, result and attribution}
A finite directed multigraph has adjacency entries $a_{uv}\in\N$, counting
arcs from $u$ to $v$, and outgoing degrees $d_v=\sum_w a_{vw}$. A
configuration is $x\in\N^V$. A legal firing at $v$ requires $x_v\geq d_v$;
it subtracts $d_v$ chips there and sends $a_{vw}$ chips to each $w$.
Firings are sequential. A legal word may stop at any configuration,
whether or not more firings are available. The empty word is allowed.
With column-vector convention, its firing matrix is
\[
 L_{wv}=a_{vw}-d_v\mathbf1_{w=v}.
\]
Thus a word of firing count $f$ has endpoint $x+Lf$. Loops cancel from
$L$ but remain in the firing threshold. All graph and configuration
integers are written in binary; multiplicities are never expanded into
lists of arcs.

Let $\reach$ consist of well-formed triples $(G,x,y)$ with $G$ nonempty,
loopless and strongly connected, and $x,y\in\N^V$, for which a legal word
from $x$ reaches $y$. Invalid encodings are rejected. Strong connectivity
and all these validity conditions are polynomial-time testable. The
one-vertex loopless case is included and has reachability exactly when
$x=y$. The number of vertices is part of the input.

\begin{theorem}\label{thm:main}
$\reach$ is coNP-complete under deterministic polynomial-time many-one
reductions.
\end{theorem}
The hardness construction uses parallel arcs, large degrees and large
chip counts. It establishes no corresponding classification for simple
graphs, unary encodings, bounded degrees or bounded total chip counts.
The membership assertion is for nonnegative configurations; it is not an
extension to every signed formulation.

The source NP-completeness theorem is Farrell--Levine
\cite[Corollary 3.2]{FL}, including its total-Laplacian hardness interface.
Hujter--Kiss--T\'othm\'er\'esz \cite[Theorem 12]{HKT} already proved the
coNP membership used here. We explain that certificate with attribution
in Section~\ref{sec:membership}; period deletion and reduced firing
counts go back to Bj\"orner--Lov\'asz \cite{BL}. The general-digraph
hardness question appears in \cite[Problem 19, arXiv v4]{HKT} and in the
Egres problem list \cite{Egres}. Its survey separates simple digraphs
from multigraphs. Hardness on our nonnegative subclass answers its
general multigraph question, including when that survey permits signed
configurations. Historical open labels do not establish openness in 2026.

There are close prior proof mechanisms. T\'othm\'er\'esz
\cite[Claim 2.4]{T} gives a deterministic recurrence-to-reachability
conversion: after one chosen legal firing from $x$, return to $x$ is
equivalent to recurrence of $x$. Proposition 2.8 there instead uses the
\emph{existence} of a recurrent representative of a nonhalting class;
it does not deterministically construct that representative as a Karp
output. The resulting Theorem 2.3 is an earlier polynomial-hierarchy
barrier. HKT's recurrent-target theorem \cite[Theorem 14]{HKT} likewise
retains a material recurrence premise. Kiss--T\'othm\'er\'esz
\cite[Section 4]{KT} already use large reserves for finite-horizon
simulation in a distance-to-nontermination/divisor-rank reduction.
Scaling and bounded-drift reserves in general are not claimed as new.
Here the specified-target argument is the gate/marker equivalence below.
These comparisons establish neither priority for that argument nor the
absence of equivalent earlier proofs.

\section{Source conventions and preprocessing}\label{sec:source}
The credited source problem asks whether a nonnegative configuration on
a strongly connected binary-adjacency directed multigraph halts
\cite{FL}. Loops are allowed. The standard halting dichotomy states that
either every maximal legal game terminates at a common stable
configuration, or every maximal game is infinite
\cite[Theorem 1]{HKT}. We use $\nh$ for valid instances in the latter
case. It is coNP-hard. To be precise about strings, a reduction to
HALTING can replace any invalid output by a fixed valid nonhalting
instance without changing its HALTING answer; complementing then gives
valid nonhalting outputs for the coNP source. Equivalently one can use
the literal complement of HALTING, including invalid strings, and send
those strings to a fixed YES reachability instance.

For valid-instance $\nh$ with invalid strings rejected, our map sends
invalid strings to the two-cycle with $X=(0,0)$ and $Y=(1,0)$, a NO
instance by conservation. For the literal-complement convention it
sends them to the same graph with $X=Y=(0,0)$, a YES instance. These
conventions are alternative total extensions, not conflicting promises.

\begin{lemma}\label{lem:preprocess}
The source can be restricted to loopless strongly connected graphs on
at least two vertices, with nonnegative configurations. Signed
intermediate configurations can also be normalized in polynomial time
without changing legal words before loop removal.
\end{lemma}
\begin{proof}
If a signed intermediate vector is used, set $k_v=\max(0,-x_v)$ and
add $k_v$ loops and $k_v$ chips at each $v$. The net firing matrix is
unchanged. After any word, its current coordinate and its threshold
both differ from their old values by $k_v$. Induction therefore gives
equivalent legality at every step, not merely equality of Laplacians.
Stability and infinite play correspond, and the new vector is
nonnegative. This supplies an explicit bridge for signed lattice
representatives in the source hardness interface.

For each vertex having $\ell_v>0$ loops, replace them by $\ell_v$ arcs
$v\to r_v$ and $\ell_v$ arcs $r_v\to v$, with a fresh relay $r_v$
initially empty. Keep other arcs and old initial counts. The old degree
is unchanged. An old word lifts by returning each relay packet
immediately after its parent fires. Conversely, in any expanded legal
word let $t_v,s_v$ be the parent and relay counts. The relay holds
$\ell_v(t_v-s_v)\geq0$, hence $s_v\leq t_v$. At each old vertex the
projected original state exceeds the expanded state by exactly that
pending return. Every projected old move is legal. An infinite
expanded word has infinitely many old moves, since relay counts are
bounded by parent counts. Thus existence of infinite play is preserved
in both directions, and the dichotomy gives the halting equivalence.
Strong connectivity survives, at most one relay is added per vertex,
and all integers retain polynomial bit length.

A singleton with $\ell$ loops never changes its count and is
nonhalting exactly when $x\geq\ell$, including $\ell=0$. Classify it
by one binary comparison and replace it by the two-cycle at $(1,0)$
for nonhalting or $(0,0)$ for halting. Thereafter $n\geq2$ and all
degrees are positive. All steps are polynomial-time.
\end{proof}

\section{The gate and marker reduction}
Let the preprocessed source have $n\geq2$ vertices, initial $x\geq0$,
total $C=\sum_v x_v$, maximum outgoing degree $D$, and incoming degrees
$i_v=\sum_u a_{uv}$. Define
\begin{equation}\label{eq:parameters}
 N=(C+1)^n,\qquad M=N+1,\qquad B=N(MD+1).
\end{equation}
The integer $N$ is a finite-state cutoff, not a polynomial running-time
bound.
\begin{lemma}\label{lem:cutoff}
The source is nonhalting if and only if it has a legal word of length $N$.
\end{lemma}
\begin{proof}
Conservation and nonnegativity put every legal source state in
$\{0,\ldots,C\}^n$, a set of at most $N$ states. A word of length $N$
visits $N+1$ states and repeats one. The intervening nonempty legal
closed word repeats forever. By the dichotomy no terminating game is
then possible. Conversely a nonhalting game supplies a length-$N$
prefix. The case $C=0$, $N=1$ is included: positive source degrees
make every vertex initially illegal.
\end{proof}

Construct $H$ with the old vertices and a gate $g$ and marker $m$:
scale every old multiplicity by $M$; add one arc $v\to g$ and $B$
arcs $g\to v$ for each old $v$; add one arc in each direction between
$g$ and $m$. The thresholds are
\[
 d^H_v=Md_v+1,\qquad d^H_g=nB+1,\qquad d^H_m=1.
\]
There are no loops, and the gate connects every old vertex and marker
in both directions. Thus $H$ is strongly connected. Put
\begin{align}
 X_v&=Mx_v+N,& X_g&=nB+1-N,& X_m&=0,\label{eq:X}\\
 Y_v&=Mx_v+MN(i_v-d_v)+B,& Y_g&=(n-1)N,&Y_m&=1.\label{eq:Y}
\end{align}
These are valid nonnegative vectors. Indeed $X_g\geq1$ and
$Y_v\geq B-MND=N$. Both totals are
$MC+(n-1)N+nB+1$, and $X\ne Y$.

\begin{lemma}\label{lem:simulation}
Before the first gate firing, the first $N$ old moves simulate the
source exactly, in every possible order. The gate cannot fire before
$N$ old moves and becomes legal after exactly $N$ such moves.
\end{lemma}
\begin{proof}
The marker remains empty before the first gate and cannot fire. After
$t$ old moves with individual counts $t_v$, let $z$ denote the source
state obtained by projecting those moves, inductively while they are
source-legal. Direct comparison of updates gives
\begin{equation}\label{eq:invariant}
 h_v=Mz_v+N-t_v,\qquad h_g=nB+1-N+t,\qquad h_m=0.
\end{equation}
Just before any of the first $N$ old moves, $t<N$ and
$0\leq t_v\leq N-1$. If $z_v\geq d_v$, then
$h_v\geq Md_v+1$; if $z_v\leq d_v-1$, then
\[
 h_v\leq M(d_v-1)+N=Md_v-1<Md_v+1.
\]
This proves legality equivalence inductively for each chosen move.
The gate formula proves the claimed threshold timing, including the
boundary just before move $N$.
\end{proof}

\begin{proposition}\label{prop:equivalence}
$(G,x)\in\nh$ if and only if $(H,X,Y)\in\reach$.
\end{proposition}
\begin{proof}
If the source is nonhalting, lift a legal word of exactly $N$ old
moves by Lemmas~\ref{lem:cutoff} and~\ref{lem:simulation}, and fire
$g$ once. Every old vertex had at least zero chips before that firing
and now has at least $B$. If its old count was $t_v\leq N$, complete
it to $N$, firing neither gate nor marker again. Every ordering of
these remaining moves is legal: ignoring all future incoming chips,
each vertex needs at most $N$ payments of $Md_v+1$, and
\[
 B=N(MD+1)\geq N(Md_v+1).
\]
Before each requested payment fewer than that vertex's allotted
payments have occurred. The final count is $f_v=N$ at old vertices,
$f_g=1$, $f_m=0$. Multiplication by $L_H$ gives exactly
\eqref{eq:Y}. This is reachability, so other available moves need not
be taken.

Conversely take \emph{any} legal word from $X$ to $Y$, with actual
counts $k$, not assumed equal to $f$. The marker balance gives
$k_g-k_m=1$, hence $k_g\geq1$. Before its first gate firing there was
no marker firing, and the gate must have received at least $N$ chips
from old moves. Therefore at least $N$ old moves preceded that first
gate. Their first $N$ moves project legally by
Lemma~\ref{lem:simulation}. Lemma~\ref{lem:cutoff} gives source
nonhalting. Delayed first gates, extra subsequent gate/marker firings
and added period vectors do not affect this argument.
\end{proof}

\begin{proof}[Polynomial-time bound]
For binary input length $s$, $n\leq s$ and the bit lengths of $C+1$
and $D$ are polynomial in $s$. The bit length of $N$ is at most
$1+n\log_2(C+1)$, with the constant retained when $C=0$; those of
$M$ and $B$ are also polynomial. Preprocessing uses at most linearly
many new vertices. The $n+2$ dimensional output is formed by
polynomially many operations on polynomial-bit integers. Integer
exponentiation in \eqref{eq:parameters} takes polynomial bit time.
The reduction writes neither a length-$N$ firing word nor unary
copies of parallel arcs. Lemma~\ref{lem:preprocess} and the total
encoding conventions complete the deterministic Karp map. Credited
source NP-completeness therefore proves coNP-hardness.
\end{proof}

\section{The credited coNP certificate}\label{sec:membership}
For completeness we give the strongly connected case of the prior
membership argument \cite[Theorem 12]{HKT}, using the period deletion
of \cite{BL}. This is an exposition of prior membership machinery.
Assume $n\geq2$. The integer kernel of $L$ is generated by a positive
primitive vector $p$. Directed spanning-tree cofactors compute $p$
up to their gcd; each has polynomial bit length. Integer-image
membership for $y-x$ and an integer solution $h$ of $Lh=y-x$ are
computable in polynomial time by integer linear algebra. If no such
solution exists, the pair is not reachable. Otherwise define
\[
 g=h-\left(\min_v\left\lfloor h_v/p_v\right\rfloor\right)p.
\]
Then $g\geq0$, $Lg=y-x$, and $g_w<p_w$ for some $w$. Every
nonnegative solution is $g+kp$ with integer $k\geq0$, so $g$ is the
unique reduced solution.

These counts need not be numerically small, but have polynomial bit
length. One can see this without a trajectory bound: fix such $w$,
delete row and column $w$ from $L$, and use its nonsingular minor
$L_w$ to write
\[
 g_{-w}=L_w^{-1}\big((y-x)_{-w}-g_wL_{-w,w}\big).
\]
The bound $g_w<p_w$ and cofactor/determinant bounds bound all
numerators in this expression polynomially in bit length; the
nonzero integer denominator has absolute value at least one.

If a legal word has count $f\geq p$, delete its first $p_v$
occurrences of each $v$. At a retained firing of $v$, deletions there
are complete and at other $u$ total $r_u\leq p_u$. Thus
\[
 (Lr)_v=-d_vp_v+\sum_u a_{uv}r_u\leq(Lp)_v=0.
\]
The retained word has at least as many chips at its firing coordinate
and remains legal. Its endpoint is unchanged. Iterating shows that
a reachable target has a legal realization with exactly count $g$.

A negative certificate, once $g$ exists, is an integer vector $u$
satisfying
\begin{equation}\label{eq:certificate}
 0\leq u\leq g,\qquad u\ne g,\qquad
 (x+Lu)_v<d_v\quad\text{whenever }u_v<g_v.
\end{equation}
It is not necessary that $u$ itself be legally realizable. In a
hypothetical legal word of count $g$, take the first firing exceeding
$u$. Its own preceding count is $u_v$ and the other counts are at
most their $u$ coordinates, so its available chips are at most
$(x+Lu)_v<d_v$, a contradiction. Conversely a maximal legal word
bounded by $g$ supplies such a proper $u$ if $g$ cannot be realized.
Its existence does not require the verifier to run that potentially
long word. The vector has polynomial bit length and all inequalities
are checked in polynomial time. Integer-image nonmembership and
invalid encodings are themselves polynomial-time decidable branches
of the complement verifier. The singleton case is decided by $x=y$.
This proves the credited coNP membership and completes
Theorem~\ref{thm:main}.

\section{Verification, source limits and disclosure}
The accompanying standard-library verification kit supplies fresh
exact finite controls for the cutoff, all-order prefix simulation,
cleanup budgets, arbitrary target counts, loop/signed interfaces,
binary sizes and the reduced-count certificate, with deliberate
faults. Checks stay active under optimized Python. They supplement
the written proof and are not universal or formal certificates.
The original proof edition distributed no executable: this package
does not claim to replay its omitted historical programs or aggregate
counts. Exact source editions, hashes, reading scopes, selected
authored audit records and individual custody exclusions are supplied;
third-party source bodies are not redistributed.

As of preparation on 15 October 2026, a bounded priority audit read
specified relevant scopes in 22 primary/related bodies, with 60
recorded queries. It established no earlier full antecedent among
those interfaces, which is not an exhaustive absence result.
Unavailable older Petri-net originals and every uninspected passage
remain unknown. Exact-result, exact-gadget and broader method priority,
independence of discovery and current worldwide openness are unknown.
The theorem is presented as a complete attributed proof answering the
historical mathematical question, without a first-discovery claim.

AI tools were used extensively in solving, checking, source comparison,
programming and writing. This note is unrefereed. Internal AI reviews
do not constitute external human peer review, journal acceptance or
proof-assistant certification. Preparation is not a publication claim.
The paper and original verification materials are licensed CC BY 4.0.

\begin{thebibliography}{9}
\bibitem{BL} A. Bj\"orner and L. Lov\'asz, Chip-firing games on directed
graphs, \emph{J. Algebraic Combin.} \textbf{1} (1992), 305--328.
Author manuscript, revised July 1992,
\url{https://www.cs.elte.hu/~lovasz/morepapers/abacus.pdf}.
\bibitem{FL} M. Farrell and L. Levine, CoEulerian graphs,
\emph{Proc. Amer. Math. Soc.} \textbf{144} (2016), 2847--2860.
Inspected arXiv v3: \url{https://arxiv.org/abs/1502.04690v3}.
\bibitem{HKT} B. Hujter, V. Kiss and L. T\'othm\'er\'esz, On the
complexity of the chip-firing reachability problem,
\emph{Proc. Amer. Math. Soc.} \textbf{145} (2017), 3343--3356.
\url{https://doi.org/10.1090/proc/13498}.
Numbering refers to arXiv v4: \url{https://arxiv.org/abs/1507.03209v4}.
\bibitem{KT} V. Kiss and L. T\'othm\'er\'esz, Chip-firing games on
Eulerian digraphs and NP-hardness of computing the rank of a divisor on
a graph, \emph{Discrete Appl. Math.} \textbf{193} (2015), 48--56.
Inspected arXiv v3: \url{https://arxiv.org/abs/1407.6958v3}.
\bibitem{T} L. T\'othm\'er\'esz, Rotor-routing reachability is easy,
chip-firing reachability is hard, \emph{European J. Combin.}
\textbf{101} (2022), 103466.
\url{https://doi.org/10.1016/j.ejc.2021.103466}.
Numbering refers to arXiv v2: \url{https://arxiv.org/abs/2102.11970v2}.
\bibitem{Egres} Egres Open, Complexity of the chip-firing reachability
problem for general digraphs; and Chip-firing (survey).
\url{https://lemon.cs.elte.hu/egres/open/Chip-firing}.
Historical returned page revisions were inspected in the source audit;
their open labels do not certify current status.
\end{thebibliography}
\end{document}
